%------------------------------------------------------------------------------%
\documentclass[a4paper,12pt,fleqn]{article}

\usepackage{calculo_varias_variaveis-1}

\ShowAnswers

%------------------------------------------------------------------------------%
\begin{document}

% A
\makeHeader{CFVVI}{Prova 2 -- Turma 3}{04/04/2026}

\bigskip
\textbf{Atenção}: a aplicação das regras do produto ou do quociente em casos
nos quais um dos termos é constante \textbf{será considerada incorreta},
por evidenciar compreensão inadequada das regras de derivação.

% 01
%------------------------------------------------------------------------------%
\exercise{25\%}
Calcule o limite
\(
    \lim_{(x,y) \to (1, -1)} \frac{(x-1)(y+1)}{(x-1)^2 + (y+1)^2},
\)
ou mostre que ele não existe

\begin{answer}
  \textbf{Atenção: os caminhos precisam passar pelo ponto $(1, -1)$}
  \bigskip

  Considerando o caminho $x=1$ temos
  \begin{align*}
    g(y)
    & = \frac{(x-1)(y+1)}{(x-1)^2 + (y+1)^2} \evalat{x=1}{} \\
    & = \frac{(1-1)(y+1)}{(1-1)^2 + (y+1)^2} \\
    & = \frac{0}{(y+1)^2} \\
    & = 0 \qquad\qquad\qquad\qquad \forall y \neq -1
  \end{align*}
  \[
    L_1
    = \lim_{y\to -1} g(y)
    = \lim_{y\to -1} 0
    = 0
  \]

  Considerando o caminho $y=x-2$ temos
  \begin{align*}
    h(x)
    & = \frac{(x-1)(y+1)}{(x-1)^2 + (y+1)^2} \evalat{y=x-2}{} \\
    & = \frac{(x-1)(x-2+1)}{(x-1)^2 + (x-2+1)^2}\\
    & = \frac{(x-1)(x-1)}{(x-1)^2 + (x-1)^2}\\
    & = \frac{(x-1)^2}{2(x-1)^2}\\
    & = \frac{1}{2} \qquad\qquad\qquad\qquad \forall x\neq 1
  \end{align*}
  \[
    L_2
    = \lim_{x\to 1} h(x)
    = \lim_{x\to 1} \frac{1}{2}
    = \frac{1}{2}
  \]
  O limite não existe, pois assume valores diferentes em cada caminho
\end{answer}

% 02
%------------------------------------------------------------------------------%
\exercise{25\%}
Calcule $\DM{f}{x}{y}$ para
\(
  f(x, y) = x^2 \ln(xy) + e^{xy} + \sin(xy),
\)
na ordem indicada pela notação

\begin{answer}
  Como
  \[
    \DM{f}{x}{y} = \D{}{x} \left(\D{f}{y}\right)
  \]
  Primeiro calculamos a derivada de $f$ por $y$
  \begin{align*}
    \D{f}{y}
    & = \D{}{y} \left( x^2 \ln(xy) + e^{xy} + \sin(xy) \right)
    \\[2mm]
    & = \D{}{y} \left( x^2 \ln(xy) \right)
      + \D{}{y} \left( e^{xy} \right)
      + \D{}{y} \left( \sin(xy) \right)
    \\[2mm]
    & = x^2 \D{}{y} \left( \ln(xy) \right)
      + e^{xy} \D{}{y} \left( xy \right)
      + \cos(xy) \D{}{y} \left( xy \right)
    \\[2mm]
    & = x^2 \frac{1}{xy} \D{}{y} \left(xy\right)
      + e^{xy} x
      + \cos(xy) x
    \\[2mm]
    & = \frac{x^2}{y} + xe^{xy} + x\cos(xy)
  \end{align*}
  Calculamos agora a derivada mista
  \begin{align*}
    \DM{f}{x}{y}
    & = \D{}{x} \left(\D{f}{y}\right)
    \\[2mm]
    & = \D{}{x} \left( \frac{x^2}{y} + xe^{xy} + x\cos(xy) \right)
    \\[2mm]
    & = \D{}{x} \left( \frac{x^2}{y} \right)
      + \D{}{x} \left( xe^{xy} \right)
      + \D{}{x} \left( x\cos(xy) \right)
    \\[2mm]
    & = \frac{1}{y} \D{x^2}{x}
      + \D{x}{x} e^{xy}
      + x \D{e^{xy}}{x}
      + \D{x}{x} \cos(xy)
      + x\D{}{x} \left( \cos(xy) \right)
    \\[2mm]
    & = \frac{2x}{y}
      + e^{xy}
      + x e^{xy}  \D{xy}{x}
      + \cos(xy)
      + x\left(-\sin(xy)\right)\D{xy}{x}
    \\[2mm]
    & = \frac{2x}{y}
      + e^{xy}
      + xy e^{xy}
      + \cos(xy)
      - xy\sin(xy)
  \end{align*}
\end{answer}

% 03
%------------------------------------------------------------------------------%
\exercise{25\%}
Use a regra da cadeia para calcular $\frac{dz}{dt}$, onde
\(
  z = \ln\left(x^2 + y^2\right),
\)
\(
  x = e^{t^2},
\)
\(
 y = \sin t
\)

\begin{answer}
\begin{multicols}{2}
  Notamos que \(z = z(x, y)\), \(x = x(t)\) e \(y = y(t)\),
  portanto pela regra da cadeia temos
  \[
    \frac{dz}{dt}
    = \D{z}{x} \frac{dx}{dt}
    + \D{z}{y} \frac{dy}{dt}
  \]
  Calculando cada derivada
  \begin{align*}
    \D{z}{x}
    & = \D{}{x} \left( \ln\left(x^2 + y^2\right) \right)
    \\[2mm]
    & = \frac{1}{x^2 + y^2} \, \D{}{x} \left( x^2 + y^2\right)
    \\[2mm]
    & = \frac{2x}{x^2 + y^2}
  \end{align*}
  \begin{align*}
    \D{z}{y}
    & = \D{}{y} \left( \ln\left(x^2 + y^2\right) \right)
    \\*[2mm]
    & = \frac{1}{x^2 + y^2} \, \D{}{y} \left( x^2 + y^2\right)
    \\*[2mm]
    & = \frac{2y}{x^2 + y^2}
  \end{align*}

  \[
    \frac{dx}{dt}
    = \frac{d e^{t^2}}{dt}
    = e^{t^2} \frac{d t^2}{dt}
    = e^{t^2} 2t
    = 2t e^{t^2}
  \]

  \[
    \frac{dy}{dt}
    = \frac{d \sin t}{dt}
    = \cos t
  \]

  \begin{align*}
    \frac{dz}{dt}
    & = \D{z}{x} \frac{dx}{dt}
      + \D{z}{y} \frac{dy}{dt}
    \\[2mm]
    & = \frac{2x}{x^2 + y^2} \, 2t e^{t^2}
      + \frac{2y}{x^2 + y^2} \, \cos t
    \\[2mm]
    & = \frac{2}{x^2 + y^2}
        \left(
          x 2t e^{t^2}
        + y \cos t
        \right)
    \\[2mm]
    & = \frac{2}{\left(e^{t^2}\right)^2 + \left(\sin t\right)^2}
        \left[
          e^{t^2} 2t e^{t^2}
          + \sin t \cos t
        \right]
    \\[2mm]
    & = \frac{2}{e^{2t^2} + \sin^2 t}
        \left[
          2t e^{2t^2}
          + \sin t \cos t
        \right]
  \end{align*}
\end{multicols}
\end{answer}

% 04
%------------------------------------------------------------------------------%
\exercise{25\%}
Use derivação implícita para avaliar \(\frac{dy}{dx},\)
sabendo que a equação
\(
  x^2 + y^2 + e^{xy} = \ln(x + y)
\)
define $y$ como função de $x$

\begin{answer}
  Vamos usar a fórmula
  \[
    \frac{dy}{dx} = - \frac{F_x}{F_y}
  \]
  com
  \[
    F(x, y) = x^2 + y^2 + e^{xy} - \ln(x + y)
  \]

  Calculando as derivadas parciais
  \begin{align*}
    F_x
    & = \D{F}{x}
    \\[2mm]
    & = \D{}{x} \left(x^2 + y^2 + e^{xy} - \ln(x + y)\right)
    \\[2mm]
    & = \D{x^2}{x}
      + \D{y^2}{x}
      + \D{e^{xy}}{x}
      - \D{}{x} \left(\ln(x + y)\right)
    \\[2mm]
    & = 2x + 0
      + e^{xy}\D{xy}{x}
      - \frac{1}{x + y}\,\D{}{x} (x + y)
    \\[2mm]
    & = 2x + ye^{xy} - (x + y)^{-1}
  \end{align*}

  \begin{align*}
    F_y
    & = \D{F}{y}
    \\[2mm]
    & = \D{}{y} \left(x^2 + y^2 + e^{xy} - \ln(x + y)\right)
    \\[2mm]
    & = \D{x^2}{y}
      + \D{y^2}{y}
      + \D{e^{xy}}{y}
      - \D{}{y} \left(\ln(x + y)\right)
    \\[2mm]
    & = 0 + 2y
      + e^{xy}\D{xy}{y}
      - \frac{1}{x + y}\,\D{}{y} (x + y)
    \\[2mm]
    & = 2y + xe^{xy} - (x + y)^{-1}
  \end{align*}
  Substituindo na fórmula temos
  \[
    \frac{dy}{dx}
    = - \frac{F_x}{F_y}
    = - \frac{2x + ye^{xy} - (x + y)^{-1}}{2y + xe^{xy} - (x + y)^{-1}}
  \]
\end{answer}

%------------------------------------------------------------------------------%
\end{document}
%------------------------------------------------------------------------------%

